% Part: methods
% Chapter: proofs
% Section: resources

\documentclass[../../../include/open-logic-section]{subfiles}

\begin{document}

\olfileid{mth}{prf}{res}

\section{আরও পাঠসামগ্রী}

গণিতে প্রমাণ করার পদ্ধতি নিয়ে অনেক উপকারী বই
আছে। বিশেষ করে \emph{How to Read and do Proofs: An
Introduction to Mathematical Thought Processes}
\citep{Solow2013} এবং \emph{How to Prove It: A
Structured Approach} \citep{Velleman2019} দেখো।
\href{https://richardhammack.github.io/BookOfProof/}{\emph{Book
of Proof}} \citep{Hammack2013} এবং
\href{https://scholarworks.gvsu.edu/books/7/}{\emph{Mathematical
Reasoning}} \citep{Sandstrum2019} প্রমাণের উপর
লেখা বই; অনলাইনে বিনা মূল্যে পাওয়া যায়।
দর্শনের শিক্ষার্থীদের কাছে গাণিতিক যুক্তিচর্চার
প্রাথমিক পাঠ হিসেবে \emph{More
Precisely: The Math you need to do Philosophy}
\citep{Steinhart2018} উপযোগী হতে পারে।
% BN-SRC-837 TARGET-CORRECTION-BEGIN
\par\noindent\textbf{পাঠ-সংশোধন।} বইগুলির শিরোনাম মূল ভাষায় রাখা হয়েছে। হ্যাম্যাকের বইয়ের
অনলাইন প্রবেশপথ এখন লেখকের নিজস্ব বই-পৃষ্ঠা; সেখানে
বর্তমান ডাউনলোড তৃতীয় সংস্করণের ৩.৪ সংশোধন, মূলপাঠে
উল্লিখিত ২০১৩ সালের সংস্করণ থেকে আলাদা। সুন্ডস্ট্রমের
বইয়ের এখানে দেওয়া বিশ্ববিদ্যালয়-পৃষ্ঠাটি ২০১৩ সালের
সংস্করণের; লেখকের নাম ও সেই পৃষ্ঠার প্রকাশকাল অনুসারে
গ্রন্থপঞ্জি মেরামত হয়েছে। ২০১৯ সালের অন্য সংস্করণ নেই—
এমন দাবি করা হচ্ছে না।
% BN-SRC-837 TARGET-CORRECTION-END

ইন্টারনেটেও প্রমাণ নিয়ে নানা সংক্ষিপ্ত নির্দেশিকা
আছে। যেমন,
\href{https://math.berkeley.edu/~hutching/teach/proofs.pdf}{``Introduction
  to Mathematical Arguments''} \citep{Hutchings2003} এবং
\href{https://eugeniacheng.com/wp-content/uploads/2017/02/cheng-proofguide.pdf}{``How to
  write proofs''} \citep{Cheng2004}।

\subsection{উৎসাহ জোগানোর ভিডিও}

বাড়ির কাজ করার উৎসাহ পাচ্ছ না? মন খারাপ?
এই ভিডিওগুলি হয়তো সাহায্য করবে!

\begin{itemize}
\item \url{https://www.youtube.com/watch?v=ZXsQAXx_ao0}
\item \url{https://www.youtube.com/watch?v=BQ4yd2W50No}
\item \url{https://www.youtube.com/watch?v=StTqXEQ2l-Y}
\end{itemize}

\end{document}
